% Armijo backtracking line search: start from an initial step along a descent
% direction and halve it until the sufficient decrease condition holds.
% Symbols follow Alg. 4.2 (backtracking) in Martins & Ning, Engineering Design
% Optimization (mdobook.github.io): initial step \alpha_init, sufficient
% decrease factor \mu_1, backtracking factor \rho.
% Cropped to its own bounding box so it can be dropped onto a Xournal++ page
% as an image.
%
%   latexmk -pdf armijo-pseudocode.tex
\documentclass[border=6pt,varwidth=11cm]{standalone}

\usepackage{algorithm2e}
\usepackage{amsmath}

\RestyleAlgo{plain}          % no float, no ruled box
\DontPrintSemicolon
\SetKwComment{Comment}{// }{}
\SetKwInOut{Input}{input}
\SetKwInOut{Output}{output}

\begin{document}
\begin{minipage}{\linewidth}
\begin{algorithm}[H]
  \textbf{Armijo Backtracking Line Search}
  \BlankLine
  \Input{current point $\vec{x}_k$, descent direction $\vec{p}_k$,
         initial step $\alpha_\text{init}$ (1 for Newton),
         sufficient decrease factor $\mu_1$ (e.g.\ $10^{-4}$),
         backtracking factor $\rho$ (e.g.\ $0.5$)}
  \Output{step size $\alpha_k$}
  \BlankLine
  $\phi(\alpha) \leftarrow f(\vec{x}_k + \alpha\, \vec{p}_k)$ \tcp*{$f$ along the line}
  $\phi'(0) \leftarrow \nabla f(\vec{x}_k)^\mathsf{T} \vec{p}_k$ \tcp*{slope at the start, $< 0$}
  $\alpha \leftarrow \alpha_\text{init}$\;
  \While{$\phi(\alpha) > \phi(0) + \mu_1\, \alpha\, \phi'(0)$}{
    $\alpha \leftarrow \rho\, \alpha$ \tcp*{shrink the step}
  }
  \Return $\alpha_k = \alpha$\;
\end{algorithm}
\end{minipage}
\end{document}
